\section*{Summary}

The rate of any channel is a linear functional of the velocity distribution of the halo, with a nonnegative rate kernel per bin that the master formula fixes at each velocity (\cref{sec:21.1}). For a response averaged over directions the kernel depends on the speed alone, and the rate reduces to an integral of the mean inverse speed $\eta(v_{\mathrm{min}})$. A threshold mass $2E_{\mathrm{th}}/v_{\mathrm{max}}^2$ is a floor, never an edge: the onset depends on the momentum transfer, the rate vanishes at the floor, and the floor moves with the Earth speed. An exposure that includes June reaches \SI{3.9}{\percent} lower for silicon.

A limit valid for every halo is a semi-infinite linear program in a nonnegative measure (\cref{sec:21.2}). For a direction-averaged response its solution is one stream, exactly even with several bins. With the pair bins of the calibrated cascade, dropping the standard halo weakens a silicon limit by at least 2.00 at \SI{10}{MeV} and 3.70 at \SI{2}{MeV}.

Near its low-mass edge a limit follows a power law (\cref{sec:21.3}). The rate rises as $(\mchi - m_{\mathrm{edge}})^{p_{\mathrm{on}} + k_{\mathrm{tail}} + 7/2}$, where $p_{\mathrm{on}}$ is the onset exponent of the response and $k_{\mathrm{tail}}$ that of the escape tail of the halo, so the free-pair edge of silicon rises with the power 13/2 in the standard halo and 8 to 10 in simulated ones. Its position depends on $E_{\mathrm{on}}/v_{\mathrm{max}}^2$ alone, so a gap shift of \SI{10}{meV} is equivalent to $-\SI{3.3}{km.s^{-1}}$ in escape speed and weakens a limit \SI{0.005}{MeV} above the edge by a factor of 231.

The theory uncertainties of each channel enter the likelihood as nuisance parameters and become a band on the limit (\cref{sec:21.4}). In the likelihood of one exposure a profiled nuisance parameter costs a limit at second order in its size, an envelope at first order. Each acts on the limit through its sensitivity $\beta_k$, the band is $\ee^{\pm\sigma_\beta}$, and profiling in the Neyman construction costs $\sigma_\beta^2\mathsf{t}'/2$ in the logarithm, exactly for one bin: a factor 1.427 for the silicon limit at \SI{3}{MeV}. For silicon with a heavy mediator the band is a factor of 3.4 at \SI{1}{MeV}, dominated by the spread of tabulations, and 1.48 at \SI{3}{MeV}, where two calculations would remove \SI{93}{\percent} of its variance.

Without a background model the maximum-gap method uses the shape of the signal (\cref{sec:21.5}). A combination of searches is one likelihood with one set of theory parameters; splitting a shared normalization parameter per exposure makes the combined limit about \SI{5}{\percent} too strong in the example.

Among the eleven survivors germanium leads the electron-recoil floor and the threshold in quanta, TeO$_2$ the heavy-mediator single-phonon rate, diamond the nuclear edge, LiF the polar rate and GaAs the polar reach (\cref{sec:21.7}). The leaders of the full pool of 42 crystals are conditional survivors. The lead salts rest on estimated pair energies and carry bulk $^{210}$Pb, which puts 724 events per kilogram, day and keV below \SI{1}{keV} for every becquerel per kilogram; CsI rests on measured elastic constants but has no demonstrated millikelvin readout.
